\chapter*{\begin{Large}PRESENTACIÓN\end{Large}}
\addcontentsline{toc}{chapter}{PRESENTACIÓN}

Lejos ya del delirio inicial que fomentó la incursión del tema de la Telemedicina en los países subdesarrollados, en donde muchos la concebían como panacea para la problemática en el cubrimiento y la calidad de los servicios de salud prestados a la población, con un mayor aplomo y madurez teórica – práctica se implementa la tercera fase de desarrollo del \textbf{Sistema Integrado de Información en Telemedicina}. Ahora es claro que la equidad en la prestación de servicios de salud va más allá del mero despliegue de redes de telecomunicaciones y debe involucrar cambios, que evidentemente se están presentando, en la legislación, las prácticas y  métodos operativos, la educación y la integración multidimensional de nodos humanos, técnicos y tecnológicos en las denominadas \textit{redes de telemedicina de tercera generación}.

En los últimos años el diseño, desarrollo e implementación de políticas innovadoras en el área de la prestación de servicios de salud – perjudiciales o no, ha logrado que las instituciones y entidades optimicen sus procesos para el manejo y la gestión eficiente de los recursos médicos, financieros y operativos. Sin embargo, las reglas que fomentan la competencia entre prestadoras de servicios de salud genera que, en busca del ideal beneficio del usuario, se segregue la condición humana en aras de una mayor rentabilidad.

El presente documento contiene los resultados obtenidos en el desarrollo del proyecto intentando abarcar con enfoque holístico la evolución conceptual, teórica, arquitectónica y normativa que ha presentado el proyecto con respecto a sus versiones anteriores. Es en la tercera fase que el \textbf{Sistema de Información en Telemedicina} transciende de la dimensión técnica y se consolida como un modelo que es útil para demostrar la importancia y validez que tienen las redes de trabajo colaborativo y el movimiento de Software Libre en el despliegue de soluciones costo efectivas en el área de la telesalud, la administración de servicios médicos y la telemedicina.


\begin{figure}[!htpb]
\begin{minipage}[r]{0.8\textwidth}
\begin{flushright}
\textbf{Lilia Edith Aparicio Pico}

Grupo de Investigación en Telemedicina

Universidad Distrital Francisco José de Caldas
\end{flushright}
\end{minipage}
\begin{minipage}[r]{0.15\textwidth}
 \includegraphics[width=20mm, height=26mm]{edith.png}
\end{minipage}
\end{figure}